Ce thème est le seul où formaliser oblige d'abord à \emph{choisir un modèle}. « On tire une
boule au hasard » n'est pas une proposition mathématique tant qu'on n'a pas dit sur quel
ensemble et avec quelle loi : la phrase décrit une expérience, pas un objet.

Au collège, le modèle est toujours le même — une probabilité uniforme sur un ensemble fini —
et tout se ramène à du dénombrement. Le lycée l'élargit deux fois : d'abord aux variables
aléatoires et à la loi binomiale, qui restent finies ; puis aux lois continues — uniforme,
exponentielle, normale — où une probabilité devient une intégrale et où le dénombrement ne
sert plus à rien.

C'est aussi le thème où l'écart entre le programme et ce qu'on sait démontrer est le plus
large, dans les deux sens. L'espérance d'une loi binomiale, admise en première, est démontrée
ici — et la bibliothèque ne la contenait pas. Le théorème de Moivre–Laplace et les intervalles
de fluctuation, admis eux aussi, ne le sont pas : ils demandent un théorème central limite
dont la forme scolaire est un travail de traduction à part entière.
